\section{Protocolo experimental}
\label{sec:protocolo}
\subsection{Diseño ejecutado y selección de las series}
Se compararon sockets TCP con JSON y gRPC con Protocol Buffers en
procesos Python sobre localhost. Los nodos de eventos son A, B y C;
las pruebas de RTT utilizan servidores de eco específicos y uno o tres
clientes. TCP proporciona un flujo fiable y ordenado, mientras HTTP/2
define los flujos y tramas del transporte utilizado por gRPC
\cite{eddy2022tcp,thomson2022http2}. El formato de serialización forma
parte de la implementación comparada \cite{jackson2024streaming}.

Se incluyeron las sesiones de Melanir del 19/11, Mourad del 22/11 y
DESKTOP-5MJCPHU del 19 y 21/11. Cada equipo tiene datos de ambos
mecanismos. Las primeras mediciones nativas de sockets del 19/11 se
conservan en la raíz de \texttt{data/}, pero no se mezclan con las
mediciones posteriores de Mourad: carecen de duración total y no son
necesarias para el conjunto final. La selección exacta se guarda en
\texttt{analysis/resultados-finales-2026-11-22/sesiones.json}.

\begin{table}[htbp]\centering\small
\caption{Condiciones ejecutadas por tamaño y mecanismo.}
\begin{tabularx}{\linewidth}{lXX}
\toprule
Condición & Secuencial & Concurrente \\
\midrule
Clientes & Uno, 100 solicitudes por corrida & Tres, 100 solicitudes cada uno \\
Solicitud serializada & 100 y 10\,000 bytes & 100 y 10\,000 bytes \\
Calentamiento & Una corrida etiquetada y excluida & Una corrida etiquetada y excluida \\
Repeticiones válidas & Cinco en todos los equipos & Cinco, salvo gRPC de Steven: una por tamaño \\
\bottomrule
\end{tabularx}\end{table}

Los tamaños ejecutados son los de la solicitud serializada completa,
no 10\,240 bytes de contenido útil. Los clientes ajustan el JSON con
su delimitador y el mensaje protobuf mediante \texttt{ByteSize()}.
Las conexiones se reutilizan durante los 100 envíos de cada cliente.
Cada equipo se mantiene separado en tablas y figuras; no se calcula
una media que combine máquinas distintas.

\subsection{Entornos registrados}
\begin{table}[htbp]\centering\small
\caption{Equipos utilizados para las mediciones finales.}
\begin{tabularx}{\linewidth}{lXXX}
\toprule
Dato & Melanie & Angel & Steven \\
\midrule
Equipo & Melanir & Mourad & DESKTOP-5MJCPHU \\
Sistema & Windows 11 Home, 26200 & Windows 11 Pro & Windows, 26300 \\
Procesador & Ryzen 7 5700U & Core i7-10700 & Core i7-13620H \\
RAM registrada & 15,35 GiB & 15,8 GiB & 15,73 GB \\
Python & 3.14.4 & 3.14.2 & 3.12.10 \\
\bottomrule
\end{tabularx}\end{table}

Los datos proceden de \texttt{data/entorno\_melanie.json},
\texttt{data/entorno\_angel.json} y \texttt{data/readme.md}; se
conservan las unidades de RAM declaradas por cada integrante.
Las versiones fijadas son grpcio y grpcio-tools 1.85.1, protobuf
7.36.2, pandas 3.1.1 y matplotlib 3.11.2. Los archivos
\texttt{sesion.json} documentan las fechas y el cliente de medición
utilizado. No se registró de forma continua el uso de CPU ni la
carga de fondo; no se presume que esas condiciones fueran iguales.

\subsection{Límites del cronometraje}
La latencia individual usa un contador de alta resolución en el cliente.
En sockets, el JSON se construye antes del cronómetro; en gRPC se
prepara el objeto antes, pero la llamada incluye la serialización interna.
Los tamaños y las cantidades coinciden, aunque las etapas medidas no
son idénticas. No se atribuye toda diferencia al transporte.

Melanie, Angel y la sesión TCP de Steven emplearon el programa común
\texttt{analysis/medir\_equipo.py}. La duración de cada corrida incluye
preparación y cierre de los clientes y conexiones y, en concurrencia,
el grupo de hilos. Una barrera coordina el trabajo antes de preparar
las conexiones. Los CSV gRPC de Steven del 19/11 proceden del cliente
original: en concurrencia la barrera se libera después de preparar los
canales, de modo que la duración excluye esa preparación. Sus tasas
se muestran con esa salvedad, no como una comparación aislada y
equivalente del transporte entre sockets y gRPC.

\subsection{Análisis y reproducción}
Los CSV identifican equipo, corrida, modo, bytes serializados, mensaje,
RTT, duración de corrida y calentamiento. La duración repetida en cada
fila se cuenta una sola vez por corrida. Se conservaron los originales
y se comprobó su integridad con \texttt{data/SHA256SUMS}.

El programa \texttt{analysis/integrar\_sesiones.py} verifica los hashes,
selecciona las series y reutiliza el analizador existente. Desde la raíz:
\begin{quote}\small\ttfamily
python analysis/integrar\_sesiones.py --salida\par
analysis/resultados-finales-revision
\end{quote}
La carpeta de salida debe ser nueva. Se calculan media, mediana,
desviación estándar muestral por mensaje y P95 mediante interpolación
lineal. El rendimiento es mensajes completados divididos entre la suma
de duraciones válidas. El gráfico de barras utiliza la media de las
medias de corrida y una desviación estándar entre esas medias.
Con una sola corrida concurrente gRPC de Steven, no se puede estimar
esa variabilidad y no se dibuja barra de error. Los intervalos de las
otras condiciones tampoco son intervalos de confianza.
